\documentclass[12pt]{article}
\usepackage{times}
\usepackage[margin=1in]{geometry}
\setlength{\headheight}{14.49998pt}
\usepackage{amsmath}
\usepackage{amssymb}
\usepackage{amsthm}
\usepackage{graphicx}
\usepackage{fancyhdr}
\usepackage{hyperref}
\usepackage{parskip}
\usepackage{sourcecodepro}
\usepackage{listings}
\usepackage{color}

% Header setup
\pagestyle{fancy}
\fancyhf{}
\rhead{Name: \underline{\hspace{3cm}}}
\lhead{APMA 4302 - Methods}
\cfoot{\thepage}

\begin{document}

\title{APMA 4302 Methods - Project Proposal}
\author{Marc Spiegelman}
\date{\today}
\maketitle

\paragraph{Proposal: background} The purpose of the project is 2-fold

\begin{itemize}
  \item To give you an opportunity to work on a problem you are truly interested in and to explore computational methods in more depth than we have time for in class.
  \item To demonstrate to me that you understand the overall workflow of problem solving from problem selection to implementation, careful testing and discussion.
\end{itemize}

\paragraph{Instructions:} In roughly a page, propose a problem that you would like to work on for the independent class project.

Your proposal should include the following:

\begin{itemize}
  \item A  statement of the general problem you want to work on
  \item A  statement of why the problem is important (to you or to a larger audience, or both)
  \item A  statement of the computational tools you expect to use or want to explore.
  \item  A  statement of what you hope to gain from the project
\end{itemize}

\paragraph{Notes/Comments:} I am relatively agnostic to the problem you choose as long as it involves high-performance/parallel computing of some flavor.
You are not limited to methods we talked about in class (e.g. you could explore spectral methods, ML algorithms etc), but it should be relavant to HPC in some way.

Some good approaches for picking a problem include
\begin{itemize}
  \item Finding  a classic paper on some problem that you would like to learn by reproducing some of those results
  \item There is some relevant technique/problem you would like to explore related to your research or some other interest you have (but you should not reproduce something you have already done in the past).
  \item There is a subject we covered that you want to explore in more depth.
\end{itemize}

The purpose of the proposal is to get you thinking about what you want to do and to give me an opportunity to provide feedback. I will be happy to discuss potential project ideas with you at any time.
\end{document}